\section{Síntesis y ampliación}

\subsection{Resumen de los puntos clave de toda la clase}
\begin{table}[H]
\centering
\begin{tabular}{p{2.5cm}p{3.7cm}p{3.8cm}p{2.6cm}}
\toprule
\textbf{Módulo de capacidad} & \textbf{Juicio central de la clase} & \textbf{Acción de ingeniería} & \textbf{Riesgo principal} \\
\midrule
Reasoning & Para actuar mejor, no para tener una cadena de pensamiento más larga & Incorporar reflexión, replanificación y verificación de acciones & Costo de tokens y amplificación de alucinaciones \\
Memory & La memoria en capas es el núcleo de la estabilidad en tareas largas & Diseñar estrategias de lectura, escritura, compresión y eliminación & Contaminación de la memoria y desajuste en la recuperación \\
Planning & La planificación dinámica es superior a una lista estática & Monitorear las desviaciones y activar la replanificación & El exceso de planificación reduce la eficiencia \\
Multi-agent & Una ganancia condicional, no una ganancia por defecto & Definir primero el protocolo, y después agregar roles & Sobrecarga de comunicación y desviación en la colaboración \\
Evaluation & Hay que observar las métricas de proceso de toda la cadena & Establecer tracing y un schema unificado & Ver solo la puntuación final produce evaluaciones erróneas \\
\bottomrule
\end{tabular}
\caption{Marco de implementación sistemática de Reasoning, Memory y Planning}
\end{table}

\begin{importantbox}{Resumen en una frase}
La siguiente etapa de competencia entre los Language Agent no consiste en quién pueda hacer el demo más vistoso, sino en quién pueda convertir el razonamiento, la memoria y la planificación en un sistema de ingeniería verificable, operable y escalable.
\end{importantbox}

\subsection{Lecturas adicionales}
\begin{itemize}
    \item Yao et al., \textit{ReAct: Synergizing Reasoning and Acting in Language Models}
    \item Shinn et al., \textit{Reflexion: Language Agents with Verbal Reinforcement Learning}
    \item Park et al., \textit{Generative Agents: Interactive Simulacra of Human Behavior}
    \item Wang et al., \textit{Voyager: An Open-Ended Embodied Agent with LLMs}
    \item Blogs técnicos oficiales de OpenAI / Anthropic sobre Agent tool-use y evaluation
\end{itemize}

\subsection{Sugerencias para la práctica futura}
\begin{itemize}
    \item Primero completar en un Agent monolítico el circuito mínimo de ``observable + reversible'', y después introducir múltiples agentes;
    \item Primero construir un verificador a nivel de tarea, y después ampliar el reasoning budget;
    \item Exponer las operaciones de memoria mediante una API explícita, para evitar la acumulación implícita de contexto;
    \item Para escenarios de producción, establecer estrategias de relevo manual y degradación, priorizando la garantía de la controlabilidad.
\end{itemize}

\end{document}
